\chapter*{General Introduction}
\addcontentsline{toc}{chapter}{General Introduction}
\label{ch:introduction}

Code optimization is a fundamental area of computer science concerned with transforming programs to improve their execution time, memory usage, and energy consumption. At the compiler level, this involves selecting and parameterizing sequences of transformations such as loop tiling, interchange, fusion, vectorization, and parallelization. Traditionally, compilers rely on hand-crafted heuristics developed by domain experts, a process that is slow, fragile across architectures, and difficult to generalize to new program patterns. The core difficulty is the \textbf{phase ordering problem}: transformations interact in non-linear ways, and a locally beneficial decision can preclude a globally superior configuration~\cite{touati2014phase}. These challenges have motivated a growing interest in automated, data-driven approaches to compiler optimization.

Reinforcement learning (\gls{rl}) offers a compelling alternative to manually designed heuristics. In \gls{rl}, an agent learns to make sequential decisions by interacting with an environment and receiving a numerical reward signal that indicates the quality of its actions~\cite{sutton2018reinforcement}, which makes \gls{rl} well suited to compiler optimization. Deep reinforcement learning employs neural networks as function approximators. It has demonstrated the ability to learn sophisticated strategies in domains  ranging from game playing to adaptive architecture-aware compiler scheduling.

The Multi-Level Intermediate Representation (\gls{mlir}), introduced by Lattner et al.~\cite{lattner2021mlir}, is a compiler infrastructure within the \gls{llvm} ecosystem that supports multiple dialects at different levels of abstraction within a single representation. Crucially, \gls{mlir} exposes high-level code structures and fine-grained transformations directly to Python, establishing it as an exceptionally flexible foundation for automated, \gls{rl}-driven optimization.

The application of \gls{rl} to \gls{mlir}-based optimization was pioneered by Bendib et al.~\cite{bendib2024reinforcement} and later extended by Tirichine et al.~\cite{tirichine2025reinforcement} to support full neural network computation graphs. While their work demonstrated the feasibility of \gls{rl}-based full-program optimization, it also revealed limitations: hardware-agnostic observations, sparse terminal reward signals, and sequential \gls{lstm} encoders that struggled with long-range dependencies. This thesis addresses these limitations by introducing a \textbf{hardware-aware observation space}, a \textbf{shaped reward function} using static code analysis, and a \textbf{transformer-based loop nest encoder}. We evaluate these improvements to measure their contribution to optimization performance.

The prior work by Tirichine et al.~\cite{tirichine2025reinforcement} established the baseline for this thesis. Their agent was evaluated against PyTorch, the PyTorch compiler, and Halide \gls{rl} across deep learning operators and full neural network models (ResNet-18, VGG, MobileNetV2). It achieved competitive speedups on element-wise operators and outperformed PyTorch by $3.3\times$ on max-pooling operators, though it remained behind on compute-intensive operations (matmul, convolution) where architecture-specialized kernels dominate. On full models, the agent reached a $2.5\times$ speedup on MobileNetV2 relative to standard PyTorch execution. These results serve as the reference point for the improvements developed in this work.

This thesis is structured into three main parts. The first part introduces the background and related work on code optimization, reinforcement learning, and the \gls{mlir} infrastructure. The second part describes the design of our extended \gls{rl} agent, including its environment, action set, and observation model. The final part presents experimental results on benchmark models, evaluates the performance gains obtained through our method, and discusses the limitations and future directions of \gls{rl}-based compilation in \gls{mlir}.
